\section{Background and Problem Setup}
\label{sec:background}

% POLISH FIX 3: Restored the implemented CCS definition (mean minus standard
% deviation of off-diagonal passage similarities) to resolve the draft mismatch.
\paragraph{RAG faithfulness.}
RAG systems retrieve passages and condition a generator on the resulting
context \citep{lewis2020rag,karpukhin2020dpr}. A generated answer is faithful
when its substantive claims are supported by the retrieved context, not merely
when the answer is true in the world. This distinction matters because a model
can answer from parametric memory even when retrieval failed, or can produce
unsupported details despite apparently relevant passages
\citep{maynez2020faithfulness,ji2023hallucination}.

\paragraph{Local retrieval similarity.}
Most retrieval metrics are local: they score each query-passage pair with
embedding similarity, cross-encoder relevance, or reranking rank. Local scores
are useful, but they do not identify answer support. Three individually
relevant passages may be redundant, contradictory, or missing a bridge fact.
One lower-scoring passage may contain the answer.

\paragraph{Context coherence and CCS.}
% POLISH FIX 3: Restored the original CCS definition used by the code and
% matched-pair CSVs, resolving the mean-only wording mismatch.
The Context Coherence Score (CCS) is a retrieval-time statistic over a
retrieved context set \(C\):
\(\ccs(C)=\operatorname{mean}_{i<j}\cos(e_i,e_j)-
\operatorname{sd}_{i<j}\cos(e_i,e_j)\), where \(e_i\) and \(e_j\) are passage
embeddings. The mean term rewards mutual topical compatibility, while the
standard-deviation penalty reduces scores for uneven sets with one coherent
sub-cluster and one outlying passage.
CCS asks a set-level question that local relevance does not: do the passages
look mutually compatible? The audit treats CCS as a diagnostic candidate. It
does not assume that CCS causes faithfulness, and the matched test in
Section~\ref{sec:matched_similarity} rejects that strong interpretation.

\paragraph{Identifiability.}
We use ``identified'' in a practical evaluation sense. A RAG claim is not
identified if another unreported factor could plausibly explain the reported
effect. A local retrieval improvement does not identify support if context-set
structure is uncontrolled. A single automatic metric does not identify an
effect size if other faithfulness observers disagree. A tuned threshold does
not identify a reusable policy if off-diagonal transfer is poor. A method does
not identify deployment value if p99 latency and indexing cost are omitted.

\paragraph{Why the usual evidence is risky.}
The risky combination is small samples, one scorer, threshold tuning on the
reporting domain, no human calibration, and partial cost reporting. Each choice
can be reasonable in isolation. Together, they let a paper report a clean
faithfulness story without separating the variables needed to support it.
